\section{Discussion}
\label{sec:discussion}

\subsection{Why a Linear Contextual Model Is Enough}

The learner does not predict a waveform or an instantaneous channel gain.
Native CCA and ACK/HARQ processing have already converted those quantities into
busy periods, interruptions, successes, and delays. For the controller, the
relevant surface is the reward difference among 24 legal recovery profiles.
An 11-dimensional linear model is a useful first approximation to that small
surface, and LinUCB adds directed exploration when online samples are limited.
The UCB ablation supports the value of context; the unresolved feature
ablations do not motivate a larger nonlinear model.

The scenario differences also fit the protocol interpretation. Under light
Poisson traffic, queues repeatedly enter and leave contention, so activation
and recent service history can distinguish useful profiles. Heavy Poisson
traffic behaves more like a persistently active set, where fixed TMC is already
strong and exploration can cost utility. The combined case mixes arrivals,
periodic occupancy, and active-set turnover; there, contextual LinUCB separates
from UCB even though no single feature-group removal is decisive. These
patterns explain the aggregate gain without suggesting that adaptation should
improve every load.

\subsection{Distributed-Learner Interpretation}

When several nodes adapt, node $k$ receives reward
$R_k(\boldsymbol q_t)$ under the joint arm vector
$\boldsymbol q_t=(q_{1,t},\ldots,q_{|\mathcal K|,t})$. This interaction may be
read as a restricted stochastic contextual bandit game. The reading is valid
when peer actions are non-anticipating within the current decision interval,
their joint policy changes little over $T_{\rm adapt}$, and updates are
asynchronous or weakly coupled. Peers may learn under those conditions because
their aggregate behavior changes slowly enough over a local update
window. Rapid synchronized co-adaptation, action-conditioned retaliation, or
jamming breaks that approximation. The interpretation supplies neither a Nash
convergence result nor an adversarial-regret guarantee.

\subsection{Limitations and Future Work}

The operating region lies between two unhelpful extremes. A fixed occupancy
process needs no adaptation; a process that changes faster than context can be
collected makes history stale. Equation~\eqref{eq:timescale} instead assumes
locally distinguishable phases with enough dwell time for learning. Separate
Poisson-rate runs establish load dependence, while the event
\texttt{dynamic-combined-4x4} case is a faster stress test: nodes join every
10 global rounds, stay active for 200 rounds, and decisions occur every 32 rounds.
Neither experiment directly estimates $T_{\rm adapt}$.

Physical validation is also incomplete. The ns-3 runs use fixed positions,
disabled log-normal shadowing, one static frequency-selective channel draw, and
only three two-second seeds. A stronger study should separate propagation and
decoding measurements (SINR, PER/BLER, MCS, coherence time) from sensing
measurements (energy-detection margin and miss/false-busy rates). MAC service
needs simultaneous-access collisions, airtime, retransmissions, and delay
tails; learning dynamics needs recovery time and performance versus
$T_{\rm adapt}/T_{\rm dwell}$. These four measurement classes would show whether a
sliding-window learner, change-point reset, or nonlinear reward model is
actually warranted.
